\chapter{1953年全国卷}

\section{解答题}

\begin{problem}
解 \(\frac{x + 1}{x - 1} + \frac{x - 1}{x + 1} = \frac{10}{3}\).
\end{problem}

\begin{solution}
由分式有意义，得 \(x\ne1\) 且 \(x\ne-1\). 在此条件下，原方程两边同乘 \(3\left(x-1\right)\left(x+1\right)\)，得 \(3\left(x+1\right)^2+3\left(x-1\right)^2=10\left(x^2-1\right)\)，即 \(6\left(x^2+1\right)=10\left(x^2-1\right)\). 整理得 \(4x^2=16\)，所以 \(x=\pm2\). 这两个值均不等于 \(\pm1\)，代回原方程也成立，故原方程的解为 \(x=\pm2\).
\end{solution}

\begin{problem}
若 \(3x^{2} + kx + 12 = 0\) 之二根相等，求 \(k\).
\end{problem}

\begin{solution}
题中方程是关于 \(x\) 的二次方程，且二次项系数 \(3\ne0\). 两根相等的充要条件是判别式为零，因此 \(\Delta=k^2-4\cdot3\cdot12=0\). 于是 \(k^2=144\)，解得 \(k=12\) 或 \(k=-12\)，即 \(k=\pm12\). 此时重根为 \(-k/6\)，确实为实数，条件满足.
\end{solution}

\begin{problem}
求 \(\begin{vmatrix}3 & -1 & 1\\ 2 & 4 & 6\\ 7 & 0 & 5 \end{vmatrix}\) 之值.
\end{problem}

\begin{solution}
按第一行展开，原行列式等于 \(3\left(4\cdot5-6\cdot0\right)-\left(-1\right)\left(2\cdot5-6\cdot7\right)+1\left(2\cdot0-4\cdot7\right)\). 因而其值为 \(60+\left(10-42\right)-28=60-32-28=0\).
\end{solution}

\begin{problem}
求 \(\log_{10}\frac{300}{7} + \log_{10}\frac{700}{3} + \log_{10}1\) 之值.
\end{problem}

\begin{solution}
三个真数均为正数，可以使用对数的乘法公式，且 \(\log_{10}1=0\). 所以原式 \(=\log_{10}\left(\frac{300}{7}\cdot\frac{700}{3}\cdot1\right)=\log_{10}10000=4\).
\end{solution}

\begin{problem}
求 \(\tan\left(870^{\circ}\right)\).
\end{problem}

\begin{solution}
因为 \(870^{\circ}=5\cdot180^{\circ}-30^{\circ}\)，而正切函数的周期为 \(180^{\circ}\)，所以 \(\tan870^{\circ}=\tan\left(-30^{\circ}\right)=-\tan30^{\circ}=-\frac{1}{\sqrt{3}}\).
\end{solution}

\begin{problem}
若 \(\cos2x=\frac{1}{2}\)，求 \(x\) 之通值.
\end{problem}

\begin{solution}
在一个周期内，满足 \(\cos\theta=\frac12\) 的角为 \(\theta=2n\pi\pm\frac\pi3\)，其中 \(n\in\Z\). 令 \(\theta=2x\)，得 \(2x=2n\pi\pm\frac\pi3\)，两边除以 \(2\)，所以 \(x=n\pi\pm\frac\pi6\)，其中 \(n\in\Z\). 该表达式已经包含余弦函数的全部周期性解.
\end{solution}

\begin{problem}
两三角形相似之条件为何？（把你所知道的都写出来）
\end{problem}

\begin{solution}
设对应顶点为 \(A\leftrightarrow A'\)、\(B\leftrightarrow B'\)、\(C\leftrightarrow C'\)，则两三角形相似的判定条件有以下三种：
\begin{enumerate}
\item 两个角分别相等，即 \(\angle A=\angle A'\)、\(\angle B=\angle B'\). 此时第三个角也相等，因为三角形内角和均为 \(180^{\circ}\)，所以由两角判定两三角形相似.
\item 一个角相等，且这个角的两边分别成比例，即 \(\angle A=\angle A'\)，且 \(AB:A'B'=AC:A'C'\). 这里相等的角是两组比例边的夹角，由两边及其夹角判定两三角形相似.
\item 三组对应边分别成比例，即 \(AB:A'B'=AC:A'C'=BC:B'C'\). 由三边判定两三角形相似.
\end{enumerate}
三者分别是两角、两边及其夹角、三边的相似判定条件，均可推出 \(\triangle ABC\sim\triangle A'B'C'\).
\end{solution}

\begin{problem}
长方体之长，宽，高为 \(12\text{ 寸}\)，\(3\text{ 寸}\)，\(4\text{ 寸}\)，求对角线之长.
\end{problem}

\begin{solution}
先取长方体底面的对角线，其长度为 \(\sqrt{12^2+3^2}=\sqrt{153}\text{ 寸}\). 该底面对角线与高 \(4\text{ 寸}\) 又构成直角三角形，因此长方体对角线 \(d\) 满足 \(d^2=12^2+3^2+4^2=144+9+16=169\). 由于长度取正值，得 \(d=13\text{ 寸}\).
\end{solution}

\begin{problem}
垂直三棱柱之高为 \(6\text{ 寸}\)，底面三边之长为 \(3\text{ 寸}\)，\(4\text{ 寸}\)，\(5\text{ 寸}\)，求体积.
\end{problem}

\begin{solution}
底面三边满足 \(3^2+4^2=5^2\)，所以底面是以 \(3\text{ 寸}\)、\(4\text{ 寸}\) 为直角边的直角三角形. 底面积为 \(S_{\text{底}}=\frac12\cdot3\cdot4=6\text{ 平方寸}\). 直棱柱体积等于底面积乘高，因此 \(V=S_{\text{底}}h=6\cdot6=36\text{ 立方寸}\).
\end{solution}

\begin{problem}
球之表面积为 \(36\pi\text{ 方寸}\)，求体积.
\end{problem}

\begin{solution}
设球半径为 \(r\text{ 寸}\)，则 \(r>0\). 由球的表面积公式，\(4\pi r^2=36\pi\)，所以 \(r^2=9\). 因为半径为正，故 \(r=3\text{ 寸}\). 于是球的体积为 \(V=\frac43\pi r^3=\frac43\pi\cdot3^3=36\pi\text{ 立方寸}\).
\end{solution}

\begin{problem}
解方程组 \(\begin{cases}
x^{2} - 2xy + 3y^{2} = 9,\\
4x^{2} - 5xy + 6y^{2} = 30.
\end{cases}\)
\end{problem}

\begin{solution}
将第一式乘以 \(10\)，第二式乘以 \(3\)，两式相减并消去常数项，得 \(-2x^2-5xy+12y^2=0\)，即 \(2x^2+5xy-12y^2=0\). 因式分解得 \(\left(x+4y\right)\left(2x-3y\right)=0\)，所以必须分两种情形讨论.

当 \(x+4y=0\) 时，\(x=-4y\). 代入第一式，得 \(16y^2+8y^2+3y^2=9\)，即 \(27y^2=9\)，所以 \(y=\pm\frac{\sqrt3}{3}\)，相应地 \(x=\mp\frac{4\sqrt3}{3}\). 得到两组解 \(\left(\frac{4\sqrt3}{3},-\frac{\sqrt3}{3}\right)\) 和 \(\left(-\frac{4\sqrt3}{3},\frac{\sqrt3}{3}\right)\).

当 \(2x-3y=0\) 时，\(x=\frac32y\). 代入第一式，得 \(\frac94y^2-3y^2+3y^2=9\)，即 \(\frac94y^2=9\)，所以 \(y=\pm2\)，相应地 \(x=\pm3\). 得到两组解 \(\left(3,2\right)\) 和 \(\left(-3,-2\right)\).

上述各组解均满足第一式，并且它们满足由两原式线性消元所得的关系 \(10\times\text{第一式}-3\times\text{第二式}=0\)，所以代回即有第二式成立. 因此原方程组的解为 \(\left(x,y\right)=\left(\frac{4\sqrt3}{3},-\frac{\sqrt3}{3}\right)\)、\(\left(-\frac{4\sqrt3}{3},\frac{\sqrt3}{3}\right)\)、\(\left(3,2\right)\)、\(\left(-3,-2\right)\).
\end{solution}

\begin{problem}
\begin{enumerate}
    \item 化简 \(\sqrt{\frac{12}{25}} + \sqrt[4]{90000} + \sqrt[6]{\frac{64}{27}}\).
    \item 求 \(\left( 2x^{3} + \frac{1}{x} \right)^{12}\) 之展开式中之常数项.
\end{enumerate}
\end{problem}

\begin{solution}
\begin{enumerate}
\item
因为 \(12=3\cdot2^2\)、\(25=5^2\)、\(90000=3^2\cdot10^4\)、\(64=2^6\)、\(27=3^3\)，且各根号取算术根，所以
\(\sqrt{\frac{12}{25}}=\frac25\sqrt3\)，\(\sqrt[4]{90000}=\sqrt[4]{3^2\cdot10^4}=10\sqrt3\)，\(\sqrt[6]{\frac{64}{27}}=\sqrt[6]{\frac{2^6}{3^3}}=\frac{2}{\sqrt3}=\frac23\sqrt3\). 因而原式 \(=\left(\frac25+10+\frac23\right)\sqrt3=\frac{166}{15}\sqrt3\).

\item 由于式中含有 \(\frac1x\)，须有 \(x\ne0\). 二项式展开的通项为 \(T_{r+1}=\mathrm C_{12}^{r}\left(2x^3\right)^{12-r}\left(\frac1x\right)^r=\mathrm C_{12}^{r}2^{12-r}x^{36-4r}\)，其中 \(r=0\)，\(1\)，\(\ldots\)，\(12\). 要成为常数项，必须有 \(36-4r=0\)，所以 \(r=9\)，且该值在允许范围内. 因此常数项为 \(\mathrm C_{12}^{9}2^{12-9}=\mathrm C_{12}^{3}2^3=220\cdot8=1760\).
\end{enumerate}
\end{solution}

\begin{problem}
锐角三角形 \(ABC\) 之三高线为 \(AD\)，\(BE\)，\(CF\)，垂心为 \(H\)，求证 \(HD\) 平分 \(\angle EDF\).
\end{problem}

\begin{solution}
由 \(AD\perp BC\)、\(BE\perp CA\)，可知 \(\angle ADB=\angle AEB=90^{\circ}\)，所以 \(A\)、\(B\)、\(D\)、\(E\) 共圆. 因同弦 \(AE\) 所对的圆周角相等，得 \(\angle ADE=\angle ABE\).

又因为 \(CF\perp AB\)、\(BE\perp AC\)，所以 \(\angle BFC=\angle BEC=90^{\circ}\)，从而 \(F\)、\(B\)、\(C\)、\(E\) 共圆. 由于 \(A\)、\(F\)、\(B\) 共线，且 \(A\)、\(E\)、\(C\) 共线，有 \(\angle ABE=\angle FBE\)、\(\angle FCE=\angle FCA\). 在圆 \(FBCE\) 中，同弦 \(FE\) 所对的圆周角相等，因此 \(\angle FBE=\angle FCE\).

再由 \(AD\perp BC\)、\(CF\perp AB\)，得 \(\angle CDA=\angle CFA=90^{\circ}\)，所以 \(C\)、\(A\)、\(F\)、\(D\) 共圆. 在圆 \(CAFD\) 中，同弦 \(FA\) 所对的圆周角相等，故 \(\angle FCA=\angle FDA\). 综合以上等式，\(\angle ADE=\angle ABE=\angle FBE=\angle FCE=\angle FCA=\angle FDA\)，所以直线 \(AD\) 平分 \(\angle EDF\).

由于三角形 \(ABC\) 为锐角三角形，垂心 \(H\) 在线段 \(AD\) 上，且从点 \(D\) 看，射线 \(DH\) 与射线 \(DA\) 重合，因此 \(HD\) 也平分 \(\angle EDF\).

\solutionfigure{\bitmapfigure[max width=4cm]{img/1953/national/q4_solution_fig1.png}}
\end{solution}

\begin{problem}
已知三角形的两个角为 \(45^{\circ}\) 及 \(60^{\circ}\)，而其夹边长 \(1\text{ 尺}\)；求最小边之长及面积.
\end{problem}

\begin{solution}
设三角形为 \(\triangle ABC\)，令 \(\angle B=45^{\circ}\)、\(\angle C=60^{\circ}\)，则夹在这两个角之间的边为 \(BC=1\text{ 尺}\)，且 \(\angle A=180^{\circ}-45^{\circ}-60^{\circ}=75^{\circ}\). 三角形中最小的角是 \(\angle B\)，所以最小边是其对边 \(AC\).

由正弦定理，\(\frac{AC}{\sin45^{\circ}}=\frac{BC}{\sin75^{\circ}}=\frac{1}{\sin75^{\circ}}\)，于是
\(AC=\frac{\sin45^{\circ}}{\sin75^{\circ}}=\frac{\frac{\sqrt2}{2}}{\sin\left(45^{\circ}+30^{\circ}\right)}=\frac{\frac{\sqrt2}{2}}{\frac{\sqrt2}{2}\cdot\frac{\sqrt3}{2}+\frac{\sqrt2}{2}\cdot\frac12}=\frac{2}{\sqrt3+1}=\sqrt3-1\text{ 尺}\).

以 \(BC\) 和 \(CA\) 为两边，夹角为 \(\angle C=60^{\circ}\)，所以三角形面积为 \(S=\frac12\cdot BC\cdot CA\cdot\sin60^{\circ}=\frac12\cdot1\cdot\left(\sqrt3-1\right)\cdot\frac{\sqrt3}{2}=\frac{3-\sqrt3}{4}\text{ 平方尺}\).
\end{solution}
